% Progress manuscript for the undergraduate thesis project.
% This file intentionally preserves samplepaper.tex as the original LNCS template.
% Evidence-status convention used throughout:
%   completed evidence  = produced by existing experiment artifacts;
%   development result  = selected on Train/Validation, not held-out evidence;
%   proposed method      = specified but not yet evaluated;
%   future evaluation   = must not be described as a completed result.

\documentclass[runningheads]{llncs}

\usepackage[T1]{fontenc}
\usepackage{amsmath}
\usepackage{amssymb}
\usepackage{graphicx}
\usepackage{url}

\begin{document}

\title{Adaptive Test-Time Compute for Small Reasoning Models:\\
A Breadth--Depth Characterization and Preliminary Controller Study}
\titlerunning{Adaptive Test-Time Compute for Small Reasoning Models}

\author{Dinh-Hong Vu\inst{1}\orcidID{0000-0001-8513-3942} \and
The-Tai Phung\inst{2} \and
Thai-Son Bui\inst{2}}

\authorrunning{Dinh-Hong Vu et al.}

\institute{Natural Language Processing and Knowledge Discovery Research Group,
Faculty of Information Technology, Ton Duc Thang University,
Ho Chi Minh City, Vietnam\\
\email{vudinhhong@tdtu.edu.vn}
\and
Faculty of Information Technology, Ton Duc Thang University,
Ho Chi Minh City, Vietnam\\
\email{\{52300059,52300058\}@student.tdtu.edu.vn}}

\maketitle

\begin{abstract}
Test-time scaling can improve language-model reasoning by increasing the
generation-token allowance of each response or by sampling multiple independent
reasoning trajectories. However, applying a fixed inference configuration to
all problems may waste computation, while adapting only one compute dimension
may overlook useful breadth--depth trade-offs. This work investigates these
trade-offs for a small open-weight reasoning model through a controlled oracle
grid over sampling breadth $k\in\{1,2,4\}$ and generation-token caps
$L\in\{512,1024,2048\}$. On 99 usable GSM8K problems with Qwen3-1.7B,
increasing $L$ produces substantial accuracy gains, whereas the marginal
benefit of additional samples decreases at larger budgets. Under the same
nominal budget, one longer trajectory is more accurate and consumes fewer
realized tokens than several heavily truncated short trajectories. A
retrospective joint oracle reaches the best observed accuracy of 83.84\% with
1,023.5 generated tokens per problem, compared with 4,246.7 tokens for the
best fixed profile, demonstrating allocation headroom but not a deployable
policy. As an initial Phase-2 development result, a leakage-controlled
depth-only logistic baseline matches fixed-$L=2048$ validation accuracy while
using 12.7\% fewer generated tokens. An 8192-token feasibility diagnostic with
Qwen3-1.7B was stopped because of prohibitive latency on the available GPU, so
subsequent development moved to a separately reported Qwen3-0.6B track. On the
same frozen 30-problem MATH-500 development sample, Qwen3-0.6B accuracy rises
from 10.0\% at $L=1024$ to 36.7\% at $L=2048$ and 50.0\% at $L=4096$, while
truncation falls from 90.0\% to 43.3\%. The $1024\rightarrow2048$ gain is
significant after Holm correction; the $2048\rightarrow4096$ comparison is
inconclusive. These findings motivate a verifier-free sequential controller
that allocates an initial token cap and introduces additional samples only when
their predicted marginal utility justifies the cost. The MATH result is
development evidence rather than held-out validation; joint-controller and
final confirmatory evaluation remain future work.

\keywords{Adaptive test-time compute \and Small reasoning models \and
Breadth--depth allocation \and Adaptive inference \and Compute-efficient
reasoning}
\end{abstract}

\section{Introduction}

Reasoning-oriented language models can often improve final-answer accuracy by
spending more computation at inference time. Two common ways of increasing
test-time compute are to extend a single reasoning trajectory and to generate
multiple independent trajectories before aggregating their answers. The first
mechanism allocates more tokens to each attempt, while the second increases the
number of attempts. Chain-of-thought prompting~\cite{wei2022cot},
self-consistency~\cite{wang2023selfconsistency}, and repeated-sampling studies
such as Brown et al.~\cite{brown2024monkeys} provide evidence that both forms
of additional computation can improve reasoning performance.

Most inference configurations nevertheless assign the same budget to every
input. This is inefficient when problem difficulty and model behavior are
heterogeneous: an easy problem may be solved with one short response, whereas
a harder problem may require either a longer trajectory or several attempts.
Recent methods therefore adapt sample counts~\cite{aggarwal2023adaptive,
wang2025penny,wan2025rasc}, token budgets~\cite{han2025token}, or more general
decoding computation~\cite{snell2025scaling,damani2025hard}. These results make
adaptive compute a well-established research direction, but they also mean
that a generic ``difficulty classifier chooses a budget'' is not, by itself,
a defensible contribution.

This study focuses on a narrower resource-constrained setting. We consider a
small open-weight reasoning model, no external verifier or process reward
model, a hard generation-token cap, and valid-only answer aggregation. We
measure every generated token consumed by a policy, including exploratory
rollouts that do not determine the final answer. Our operational variables are
sampling breadth $k$, the number of independent trajectories, and token cap
$L$, the \texttt{max\_new\_tokens} limit for each trajectory. We use the word
``depth'' only as shorthand for this operational cap; it is not a claim about
the semantic depth of reasoning.

The initial experiment characterizes the full $3\times3$ grid
$k\in\{1,2,4\}$ and $L\in\{512,1024,2048\}$ for Qwen3-1.7B on GSM8K.
The characterization reveals a pronounced asymmetry. Increasing $L$ improves
accuracy significantly at every $k$, while the benefit of increasing $k$ from
2 to 4 becomes small or inconclusive at the larger token caps. Moreover,
profiles with the same nominal product $kL$ have substantially different
accuracy and realized-token cost. These observations motivate, but do not
prove, the working hypothesis \emph{depth first, breadth selectively}.

The main research question for the confirmatory study is:

\begin{quote}
For small open-weight reasoning models, can a verifier-free lightweight
sequential policy that allocates generation-token caps and selectively adds
independent samples improve the held-out accuracy--actual-generated-token
Pareto frontier over fixed profiles and depth-only or breadth-only adaptive
baselines?
\end{quote}

Two supporting questions ask when selective breadth provides utility beyond
adaptive depth and how stable the observed decisions and gains are across
reasoning datasets and a second small-model family. At the current project
stage, the Qwen3-1.7B GSM8K characterization and development-only depth
baseline remain completed historical evidence. A separate 30-problem
Qwen3-0.6B MATH-500 depth pilot has also been completed. The joint controller,
larger cross-dataset study, and final held-out evaluation are explicitly future
work.

The current contributions are:

\begin{enumerate}
  \item a controlled breadth--token-cap characterization of Qwen3-1.7B on
  GSM8K with realized-token, latency, validity, and truncation accounting;
  \item an iso-budget analysis showing that equal nominal budgets can produce
  substantially different accuracy and realized cost;
  \item retrospective fixed, depth-only, breadth-only, and joint oracle
  frontiers that quantify opportunity without presenting the oracle as a
  deployable policy;
  \item a leakage-controlled controller-development pipeline and a preliminary
  depth-only logistic baseline;
  \item model-separated Qwen3-1.7B and Qwen3-0.6B MATH-500 development pilots
  with a versioned exact-answer parser and depth-cap sensitivity analysis; and
  \item a pre-specified sequential joint-controller design and a stronger
  confirmatory evaluation protocol for subsequent work.
\end{enumerate}

\section{Related Work}

\subsection{Test-Time Scaling and Self-Consistency}

Chain-of-thought prompting elicits intermediate reasoning before a final
answer~\cite{wei2022cot}. Self-consistency improves this approach by sampling
multiple reasoning paths and selecting the most frequent answer
~\cite{wang2023selfconsistency}. Repeated-sampling studies show that solution
coverage can continue to grow with the number of attempts, although converting
coverage into final-answer accuracy may require a reliable verifier or
aggregator~\cite{brown2024monkeys}. These methods establish breadth as an
important test-time compute axis, but a fixed sample count spends the same
resources on easy and difficult inputs.

\subsection{Adaptive Sampling}

Adaptive-Consistency dynamically stops sampling as answer agreement
stabilizes~\cite{aggarwal2023adaptive}. Difficulty-Adaptive Self-Consistency
combines prior estimates of question difficulty with posterior sample
evidence~\cite{wang2025penny}. Reasoning-Aware Self-Consistency also evaluates
reasoning paths when deciding whether to stop and how to aggregate
answers~\cite{wan2025rasc}. These works are direct precedents for adaptive
breadth and must be treated as strong baselines rather than omitted from the
novelty comparison.

\subsection{Adaptive Token Budgets and Compute Allocation}

Token-Budget-Aware LLM Reasoning assigns different reasoning budgets according
to estimated problem complexity~\cite{han2025token}. Thought Calibration uses
probes over hidden representations to determine when additional thinking is
unlikely to help~\cite{wu2025thought}. Snell et al.~\cite{snell2025scaling} and
Damani et al.~\cite{damani2025hard} study more general input-adaptive
allocation of inference computation, often using learned reward signals,
hidden states, or larger training collections. A closely related 2026 preprint
formalizes average-budget allocation with a Lagrangian oracle and learns oracle
actions from cheap input features, although its experiments vary sample count
rather than jointly varying sample count and token cap~\cite{zhai2026adaptive}.

\subsection{Position of This Study}

The proposed setting is intentionally narrower than general test-time search
or verifier-guided scaling. It asks whether an external, black-box,
verifier-free policy can coordinate a hard token cap and selective resampling
for a small open-weight model, while treating realized token use as the main
cost rather than relying only on nominal budgets~\cite{wang2024tokeneconomies}.
The scientific outcome is the held-out
accuracy--realized-cost comparison against strong single-axis baselines, not
profile-classification accuracy alone. We therefore avoid claiming the first
adaptive compute allocator and instead evaluate a constrained joint allocation
question that remains insufficiently characterized for small-model,
low-overhead inference.

\section{Problem Formulation}

Let $x_i$ be a reasoning problem with ground-truth answer $y_i$. An inference
profile is
\begin{equation}
  p=(k,L)\in\mathcal{P},
\end{equation}
where $k$ is the number of independent trajectories and $L$ is the maximum
number of newly generated tokens per trajectory. The nominal upper bound is
$B(p)=kL$, while the realized generated-token cost for problem $i$ is
\begin{equation}
  C_i(p)=\sum_{j=1}^{k} t_{ij}(L),
\end{equation}
where $t_{ij}(L)\leq L$ is the actual number of tokens produced by trajectory
$j$. Nominal budget and realized cost are reported separately.

Each trajectory is parsed using a strict final-answer extractor. Let
$V_i(p)$ be the multiset of valid normalized answers among the $k$
trajectories. If $V_i(p)$ is non-empty, the system returns its majority answer,
with a deterministic lowest-sample-index rule for ties. Invalid answers do not
vote. If no valid answer exists, the profile is incorrect. We denote final
correctness by $A_i(p)\in\{0,1\}$.

The retrospective cheapest-correct joint oracle is
\begin{equation}
  p_i^{*}=\arg\min_{p\in\mathcal{P}:A_i(p)=1}
  \left(C_i(p),T_i(p),B(p),k,L\right),
\end{equation}
where $T_i(p)$ is latency and the tuple specifies the deterministic
tie-breaking order. This oracle uses ground-truth outcomes after inference. It
is an upper bound on allocation opportunity and cannot be deployed.

For a deployable controller in state $s$, a cost-sensitive action objective can
be written as
\begin{equation}
 U_{\lambda}(a\mid s)=
 \Pr(\text{correct after }a\mid s)
 -\lambda\,\mathbb{E}[\Delta C\mid s,a].
\end{equation}
The intended policy chooses an initial $L$, observes the first paid trajectory,
and decides whether the expected marginal utility justifies additional
samples. Ground truth, oracle profile labels, and ungenerated trajectories are
forbidden test-time signals.

\section{Methodology}

\subsection{Oracle Grid Characterization}

The completed characterization evaluates all nine profiles in
$\{1,2,4\}\times\{512,1024,2048\}$. Samples are nested within a token cap:
the $k=1$ profile uses sample 0, $k=2$ uses samples 0--1, and $k=4$ uses samples
0--3. Consequently, each problem requires twelve unique generations rather
than independently regenerating every profile. The sample seed is derived from
the global seed, problem identifier, repetition identifier, and sample
identifier; it is independent of both $k$ and $L$. This construction supports
paired comparisons while preventing $k$ profiles from receiving unrelated
sample pools.

\subsection{Answer Parsing and Aggregation}

The parser accepts explicit numeric final-answer formats, including GSM8K
markers, boxed answers, and explicit ``Final answer'' fields. It does not fall
back to the last number in free-form reasoning. A response may be categorized
as truncated during thinking, truncated after thinking, completed but
unparseable, or completed with a valid answer. Valid-only majority voting is
used for $k>1$. Agreement at $k=1$ is not interpreted as consensus.

One selected GSM8K item was marked for manual review. It remains in raw
artifacts for provenance but is excluded from accuracy, Pareto, and oracle
metrics. Original benchmark labels are never rewritten.

\subsection{Cost and Oracle Analysis}

Allocated budget, actual generated tokens, and latency are retained as distinct
metrics. The primary frontier is constrained by actual generated tokens because
models frequently stop before $L$, and different profiles with equal $kL$ can
consume different realized compute. Fixed strategies use the same profile for
every problem. Depth-only oracles hold $k$ fixed and choose $L$ retrospectively;
breadth-only oracles hold $L$ fixed and choose $k$; the joint oracle may choose
any profile. All restricted anchors are reported before selecting the best one.

\subsection{Leakage-Controlled Depth Baseline}

The implemented depth-only development baseline uses only features available
before generation: question and prompt length, counts of numerical quantities,
operator indicators, fraction and percentage indicators, and shallow lexical
features. The feature schema explicitly allowlists columns prefixed with
\texttt{feature\_}; correctness, observed outcomes, oracle profiles, and target
labels are forbidden model inputs.

The training target is the smallest $L$ whose observed $k=1$ response is
correct. If no observed $k=1$ response is correct, the target falls back to
$L=2048$ and retains an explicit unsolved-at-maximum-depth status. A
standardized multinomial logistic regression is selected over
$C\in\{0.1,1,10\}$, class weights in \{none, balanced\}, and ordered
decision quantiles in $\{0.5,0.7,0.9\}$. Selection first maximizes validation
policy accuracy, then minimizes actual generated tokens, followed by balanced
classification metrics and deterministic simplicity tie-breakers.

\subsection{Proposed Sequential Joint Controller}

The joint controller has not yet been implemented or evaluated. Its proposed
execution is:
\begin{enumerate}
  \item choose the initial token cap from input-only features;
  \item generate and parse sample 1;
  \item stop at $k=1$ or acquire sample 2 using input and already-observed
  validity, truncation, output-length, and cap-utilization signals;
  \item after two samples, stop at $k=2$ or acquire samples 3--4 using the
  preceding signals plus answer agreement and valid-answer fraction; and
  \item aggregate only the samples that the deployed policy paid to generate.
\end{enumerate}

The primary proposed learner is regularized logistic regression, with a shallow
gradient-boosted tree as an ablation only if sufficient minority examples are
available. A flat nine-class profile classifier is not the primary design
because the observed cheapest-correct labels are highly imbalanced and some
classes are empty.

\section{Experimental Setup}

\subsection{Model and Dataset}

The completed evidence contains two model-specific tracks. The historical
breadth--depth characterization uses Qwen3-1.7B~\cite{qwen3report}, revision
\texttt{70d244cc86ccca08cf5af4e1e306ecf908b1ad5e}, and the GSM8K test
split~\cite{cobbe2021gsm8k}. Thinking mode is enabled. Generation uses
temperature 0.7, top-$p$ 0.95, CUDA, and bfloat16. The fixed dataset manifest
contains 100 problems selected with global seed 42. After the manual-audit
exclusion, 99 problems contribute to main metrics. This grid uses one
stochastic repetition and contains 1,200 unique generation records.

Current characterization and controller development use Qwen3-0.6B, revision
\texttt{c1899de289a04d12100db370d81485cdf75e47ca}. The change is a versioned
compute-feasibility amendment after a Qwen3-1.7B 8192-token diagnostic proved
too slow on the available RTX 4060. It does not relabel or pool the 1.7B
evidence. Model-dependent oracle actions, controller targets, and latency
estimates must be regenerated for 0.6B.

The completed 0.6B development pilot uses the same decoding settings on 30
deterministically selected MATH-500 problems~\cite{hendrycks2021math}. The
frozen sample is stratified across all seven subjects and five difficulty
levels. It evaluates one trajectory per problem at
$L\in\{1024,2048,4096\}$, yielding 90 unique generation records with no
generation errors. The sample validates the smaller-model adapter,
exact-answer parser, and depth-cap range; it is not an untouched MATH-500 test
partition. All completed grids use one stochastic repetition, so current
uncertainty estimates describe variation over problems but not full
generation-seed variability.

\subsection{Phase-2 Development Split}

The existing Phase-2 manifest freezes 79 Train, 20 Validation, and 100 new
held-out GSM8K problems. Train and Validation come from the 99 usable
characterization problems; the Test problems are disjoint. The manifest stores
problem identifiers, source indices, question hashes, and ground-truth hashes.
Loading rejects cross-partition overlap or record changes. The 20-problem
Validation partition is used for model selection and is not claimed to be an
untouched final evaluation. The 100-problem held-out Test partition has not
been loaded or evaluated.

The existing controller labels and fitted depth baseline were produced with
Qwen3-1.7B. The frozen problem identifiers may be reused for the 0.6B track,
but outcomes, labels, thresholds, and fitted coefficients are not transferable
and must be regenerated before any 0.6B controller evaluation.

\subsection{Metrics}

Primary outcomes are final-answer accuracy and actual generated tokens. We also
report nominal $kL$, latency, truncation rate, valid-answer fraction, and
tokens per correct answer where appropriate. Every trajectory acquired by a
sequential policy must contribute to its cost, including exploratory or
intermediate trajectories.

\subsection{Statistical Analysis}

Profile accuracy is reported with Wilson 95\% confidence intervals. Planned
depth, breadth, and equal-budget accuracy comparisons use exact paired McNemar
tests with Holm correction within each pre-specified comparison family.
Problem-level paired bootstrap intervals summarize oracle-strategy differences.
An inconclusive test is not interpreted as equivalence. Because only one
generation repetition has been completed, problem-sampling uncertainty is
reported separately from the still-unmeasured generation-seed uncertainty.

\subsection{Baselines and Evidence Status}

The final policy evaluation is intended to include fixed profiles, a learned
depth-only controller, an agreement-stopping breadth-only controller, a simple
depth-first heuristic, the learned joint controller, and retrospective fixed,
depth-only, breadth-only, and joint oracles. The depth-only controller at
comparable realized cost is the primary baseline; the joint controller must not
be evaluated only against the most expensive fixed profile.

Table~\ref{tab:status} separates completed evidence from planned work.

\begin{table}
\caption{Evidence status at the current project stage.}
\label{tab:status}
\centering
\begin{tabular}{p{5.0cm}|p{6.0cm}}
\hline
Component & Status \\
\hline
Qwen3-1.7B/GSM8K $3\times3$ grid & Completed historical evidence \\
Cost, failure, statistical, and oracle analyses & Completed evidence \\
79/20/100 Phase-2 manifest & Completed engineering artifact \\
Qwen3-1.7B depth-only logistic controller & Historical Train/Validation result \\
Qwen3-1.7B MATH-500 depth pilot & Completed historical evidence \\
Qwen3-0.6B MATH-500 depth pilot & Completed current development evidence \\
Qwen3-0.6B GSM8K labels and depth controller & Not yet regenerated \\
Sequential joint controller & Proposed method; not yet implemented \\
100-problem GSM8K held-out test & Frozen; not evaluated \\
Larger MATH and secondary-model study & Future confirmatory evaluation \\
\hline
\end{tabular}
\end{table}

\section{Preliminary Results}

\subsection{Breadth--Depth Accuracy Grid}

Table~\ref{tab:grid} presents final-answer accuracy for all nine profiles. At
every breadth, accuracy rises sharply as $L$ increases. All six planned adjacent
depth comparisons are significant after Holm correction. Four of six adjacent
breadth comparisons are significant; the $k=2$ to $k=4$ comparisons at
$L=1024$ and $L=2048$ are inconclusive, with the latter adding only 2.02
percentage points.

\begin{table}
\caption{Final-answer accuracy on 99 evaluated GSM8K problems.}
\label{tab:grid}
\centering
\begin{tabular}{c|ccc}
\hline
$k$ & $L=512$ & $L=1024$ & $L=2048$ \\
\hline
1 & 9.09\% & 53.54\% & 70.71\% \\
2 & 17.17\% & 65.66\% & 81.82\% \\
4 & 28.28\% & 71.72\% & 83.84\% \\
\hline
\end{tabular}
\end{table}

At $L=512$, sample-level truncation is approximately 85--88\%, explaining why
additional short samples frequently replicate incomplete behavior rather than
providing useful independent solutions. Increasing $k$ from 1 to 2 is useful
at all three caps, but the marginal return from 2 to 4 decreases as $L$ grows.

\subsection{Equal Nominal Budget Does Not Imply Equal Behavior}

Table~\ref{tab:isobudget} compares three profiles with the same nominal budget
$kL=2048$. Contrary to treating tokens and samples as interchangeable, the
single long trajectory is both more accurate and less expensive in realized
tokens. The accuracy differences between $4\times512$ and each longer-depth
profile are significant after Holm correction; the difference between
$2\times1024$ and $1\times2048$ is inconclusive.

\begin{table}
\caption{Profiles with the same nominal budget $kL=2048$.}
\label{tab:isobudget}
\centering
\begin{tabular}{c|r|r|r}
\hline
Profile & Accuracy & Avg. actual tokens & Sample truncation \\
\hline
$4\times512$ & 28.28\% & 1,988.9 & 85.4\% \\
$2\times1024$ & 65.66\% & 1,616.0 & 35.9\% \\
$1\times2048$ & 70.71\% & 1,079.1 & 18.2\% \\
\hline
\end{tabular}
\end{table}

This is an empirical observation for the tested setting. It does not establish
that depth is universally superior to breadth.

\subsection{Per-Problem Oracle Preferences}

The joint cheapest-correct oracle solves 83 of 99 evaluated problems. Among
these 83, 70 select a $k=1$ profile and 13 require $k>1$ under the observed
sample pool. Fifty-five select $L=2048$, while no solved problem selects a
profile with $L=512$ as cheapest-correct. Table~\ref{tab:preferences} gives the
full distribution.

\begin{table}
\caption{Cheapest-correct joint-oracle profile distribution.}
\label{tab:preferences}
\centering
\begin{tabular}{c|r@{\qquad}c|r}
\hline
Profile & Count & Profile & Count \\
\hline
$1\times512$  & 0  & $2\times512$  & 0 \\
$1\times1024$ & 25 & $2\times1024$ & 2 \\
$1\times2048$ & 45 & $2\times2048$ & 9 \\
$4\times512$  & 0  & $4\times1024$ & 1 \\
$4\times2048$ & 1  & Unsolved       & 16 \\
\hline
\end{tabular}
\end{table}

The imbalance makes direct nine-class profile prediction a poor primary
learning formulation. It also motivates an explicit minimum number of positive
escalation examples before learning separate breadth decisions.

\subsection{Retrospective Oracle Headroom}

Table~\ref{tab:oracle} compares fixed and retrospective adaptive strategy
anchors. The best fixed profile, $4\times2048$, solves 83 problems at an
average cost of 4,246.7 generated tokens. The joint oracle reaches the same
accuracy at 1,023.5 tokens. At approximately this realized cost, the exact
budget-constrained frontier solves 83 problems for joint allocation, compared
with 70 for depth-only, 69 for breadth-only, and 53 for fixed allocation. Thus,
the joint oracle has 13.13 percentage points of compute-matched headroom over
the depth-only oracle. This is hindsight opportunity, not evidence that a
controller can identify the oracle action.

\begin{table}
\caption{Illustrative fixed and retrospective oracle anchors.}
\label{tab:oracle}
\centering
\begin{tabular}{l|l|r|r|r}
\hline
Strategy & Anchor & Solved & Accuracy & Avg. actual tokens \\
\hline
Fixed & $4\times2048$ & 83 & 83.84\% & 4,246.7 \\
Breadth-only oracle & $L=2048$ & 83 & 83.84\% & 1,270.6 \\
Depth-only oracle & $k=1$ & 70 & 70.71\% & 750.5 \\
Depth-only oracle & $k=4$ & 83 & 83.84\% & 3,118.1 \\
Joint oracle & all profiles & 83 & 83.84\% & 1,023.5 \\
\hline
\end{tabular}
\end{table}

\subsection{Preliminary Depth-Only Controller}

The selected depth-only candidate uses $C=1$, balanced class weights, and the
0.9 ordered target-depth quantile. Table~\ref{tab:depthcontroller} reports its
20-problem Validation result. The controller matches the fixed-$L=2048$
accuracy of 13/20 while reducing average actual generated tokens from 1,213.05
to 1,059.45, a 12.7\% reduction. Exact depth-target classification accuracy is
50\%.

\begin{table}
\caption{Depth-only development result on 20 Validation problems.}
\label{tab:depthcontroller}
\centering
\begin{tabular}{l|r|r|r}
\hline
Policy & Correct & Accuracy & Avg. actual tokens \\
\hline
Fixed $L=512$  & 2/20  & 10.0\% & 499.00 \\
Fixed $L=1024$ & 9/20  & 45.0\% & 843.65 \\
Fixed $L=2048$ & 13/20 & 65.0\% & 1,213.05 \\
Selected depth logistic & 13/20 & 65.0\% & 1,059.45 \\
\hline
\end{tabular}
\end{table}

This result is optimistic because the candidate and its thresholds were selected
on the same Validation set. It is neither held-out evidence nor confirmation of
the thesis claim. The frozen Test partition remains unused.

\subsection{MATH-500 Depth-Cap Development Pilots}

Table~\ref{tab:mathpilot} reports the completed Qwen3-0.6B pilot on the frozen
30-problem MATH-500 development sample. Increasing the cap produces a monotonic
accuracy gain and a corresponding reduction in truncation. The paired
$1024\rightarrow2048$ comparison contains eight improvements and no
regressions, for a 26.7-point accuracy increase (exact McNemar
$p=0.0078125$, Holm-adjusted $p=0.015625$ within the planned depth family).
The $2048\rightarrow4096$ comparison adds four correct answers and no
regressions, but is inconclusive after correction ($p=0.125$). This result does
not establish equivalence. Thirteen of 30 trajectories still hit the
4096-token cap, so the upper end of the 0.6B MATH action space remains
truncation-limited.

\begin{table}
\caption{Qwen3-0.6B MATH-500 depth pilot on 30 development problems.}
\label{tab:mathpilot}
\centering
\begin{tabular}{r|r|r|r|r}
\hline
$L$ & Correct & Accuracy & Avg. actual tokens & Truncation \\
\hline
1024 & 3/30  & 10.0\% & 984.9   & 90.0\% \\
2048 & 11/30 & 36.7\% & 1,764.4 & 60.0\% \\
4096 & 15/30 & 50.0\% & 2,804.0 & 43.3\% \\
\hline
\end{tabular}
\end{table}

The earlier Qwen3-1.7B stratum on the same frozen problems achieved 6.7\%,
26.7\%, and 53.3\% accuracy, with 93.3\%, 73.3\%, and 46.7\% truncation at the
three caps. These model tracks are descriptive and separate: 0.6B performs
better at the two lower caps but is 3.3 points lower at 4096, and one
stochastic repetition cannot support a model-superiority claim.

Manual auditing of the historical run exposed two deterministic
reference-normalization defects: a numeric base suffix and a comma-separated
numeric list. Parser v2 corrected these cases and reprocessed the same 90 raw
responses without new inference. The 0.6B pilot used parser v2 from the start.
This versioning prevents parser behavior from being confounded with model
behavior.

\section{Discussion}

\subsection{How the Evidence Supports the Design Hypothesis}

The observed grid supports allocating enough room for a usable trajectory before
spending heavily on additional samples. However, breadth should not be removed:
13 of the 83 oracle-solved problems require $k>1$, and the budget-constrained
joint oracle retains substantial headroom over depth-only allocation. The
academically safe interpretation is therefore that \emph{depth first, breadth
selectively} is an empirical observation and a hypothesis for confirmatory
testing, not a universal design principle.

The two MATH-500 model strata strengthen the narrower claim that an
insufficient token cap can dominate observed performance within the Qwen3
family. They do not establish cross-family or cross-dataset controller
generalization because the sample is development-only and evaluates depth
rather than adaptive breadth.

\subsection{Oracle Opportunity Versus Learnability}

The oracle observes correctness after all profiles have been evaluated. A
controller observes only cheap input features and signals from already-paid
trajectories. The central research risk is that the 13 breadth-dependent cases
may be too rare or too stochastic to predict. Consequently, controller success
must be defined by the held-out policy frontier rather than by the size of the
oracle gap or by profile-label accuracy.

\subsection{System Contribution}

The accompanying software separates dataset and model adapters, strict answer
parsing, inference execution, nested sample reuse, append-only generation
caching, checkpoint/resume, provenance manifests, cost accounting, controller
data construction, and statistical reporting. The ordinary non-GPU test suite
currently contains 134 passing tests. This reproducible pipeline is part of the
CS/IT contribution, but infrastructure such as LangGraph, RAG, multi-agent
execution, or a deployment platform is not required by the research question.

\section{Limitations}

The completed evidence has several important limitations:
\begin{itemize}
  \item it uses two sizes from one model family and one generation repetition;
  the MATH-500 evidence is a 30-problem development pilot rather than a held-out
  benchmark result;
  \item only 99 problems contribute to main characterization metrics;
  \item $L$ is an operational token cap and must not be interpreted as true
  reasoning depth;
  \item the retrospective oracle uses ground-truth outcomes and is not
  deployable;
  \item the depth-controller result is selected on only 20 Validation problems;
  \item oracle labels are model-, prompt-, parser-, and seed-dependent;
  \item the completed Qwen3-1.7B controller labels and coefficients cannot be
  transferred to Qwen3-0.6B;
  \item latency is conditional on the recorded hardware and execution setup;
  \item the joint controller and final held-out evaluation have not been run;
  and
  \item no controller-level cross-dataset or cross-model generalization claim
  is yet supported.
\end{itemize}

\section{Planned Confirmatory Evaluation}

The revised protocol expands the completed MATH-500 pilot without turning the
thesis into a large benchmark study. The minimum dataset suite is GSM8K plus
MATH~\cite{hendrycks2021math}. The completed Qwen3-1.7B experiments remain
model-specific historical evidence. After the 8192-token diagnostic showed
prohibitive latency on the available RTX 4060, characterization and controller
development moved to Qwen3-0.6B, pinned at revision
\texttt{c1899de289a04d12100db370d81485cdf75e47ca}. The 30-problem 0.6B
MATH-500 pilot is now complete. This is a versioned compute-feasibility
amendment, not a relabelling of the 1.7B results. Metrics from the two models
are reported separately, and a universal cross-model controller remains
outside scope.

The target development scale is at least 300 oracle-characterized problems,
with at least 100 from the second dataset and at least 30 positive examples for
each learned escalation decision. The existing Qwen3-1.7B labels cannot count
as 0.6B controller targets; GSM8K and MATH outcomes required by the selected
0.6B action space must be regenerated before fitting the controller.

For Qwen3-0.6B on MATH-500, 90.0\%, 60.0\%, and 43.3\% of trajectories
truncate at $L=1024$, 2048, and 4096. The current development action space
therefore retains $L\in\{1024,2048,4096\}$: 4096 provides directional accuracy
headroom, although its comparison with 2048 is inconclusive at this sample
size. A bounded Qwen3-1.7B $L=8192$ extension on the same sample was stopped
after two trajectories required 11.63 and 27.10 minutes and still truncated
before final answers. That partial run is retained only as feasibility evidence
and excluded from accuracy estimates. The 0.6B protocol does not include 8192.

The final baseline set will include all relevant fixed profiles, a
compute-matched input heuristic, learned depth-only allocation,
Adaptive-Consistency-style breadth stopping, a simple depth-first heuristic,
the learned joint policy, and restricted/joint oracle upper bounds. Before
opening final Test outcomes, the project will freeze the dataset manifests,
feature contract, controller, utility trade-off, and decision thresholds.

The joint method will proceed as the primary method on Validation if it either
improves accuracy over depth-only by at least three percentage points at
realized-token cost within $\pm10\%$, or reduces actual generated tokens by at
least 10\% with no more than a two-point accuracy decrease. If joint allocation
is Pareto-dominated by depth-only, or breadth is too rare to provide positive
net utility, the thesis will pivot to adaptive depth and report the failure to
convert joint-oracle headroom into a learnable policy. If minority events are
insufficient, the response will be to collect more oracle data or simplify the
decision rather than increase controller complexity.

\section{Conclusion}

The completed characterization shows that the placement of test-time compute
matters: equal nominal budgets do not produce equal accuracy or realized cost,
and short, highly truncated trajectories can make additional breadth
inefficient. Most oracle-solved problems favor $k=1$, yet a minority requires
additional samples and creates joint-allocation headroom. A preliminary
depth-only controller provides encouraging development evidence by matching
the strongest fixed-depth Validation accuracy with 12.7\% fewer generated
tokens. The current Qwen3-0.6B MATH-500 development pilot independently shows
that increasing the token cap from 1024 to 4096 raises accuracy from 10.0\% to
50.0\% while reducing truncation from 90.0\% to 43.3\%. The significant gain
from 1024 to 2048 and the inconclusive 2048-to-4096 comparison support retaining
all three caps for further characterization without claiming that 4096 is
decisively superior. These findings justify a sequential, selective-breadth
hypothesis, but they do not yet establish a deployable joint policy or
controller generalization across datasets. The next scientific step is to
regenerate 0.6B controller labels, implement the staged decisions, and evaluate
the frozen policy on unseen problems.

\begin{credits}
\subsubsection{Acknowledgements}
The authors thank the Faculty of Information Technology at Ton Duc Thang
University for supervision and research support.

\subsubsection{Disclosure of Interests}
The authors have no competing interests to declare that are relevant to the
content of this article.
\end{credits}

\begin{thebibliography}{15}

\bibitem{wei2022cot}
Wei, J., Wang, X., Schuurmans, D., Bosma, M., Xia, F., Chi, E.H., Le, Q.V.,
Zhou, D.: Chain-of-thought prompting elicits reasoning in large language
models. In: Advances in Neural Information Processing Systems 35 (2022)

\bibitem{wang2023selfconsistency}
Wang, X., Wei, J., Schuurmans, D., Le, Q.V., Chi, E.H., Narang, S.,
Chowdhery, A., Zhou, D.: Self-consistency improves chain of thought reasoning
in language models. In: International Conference on Learning Representations
(2023)

\bibitem{brown2024monkeys}
Brown, B., Juravsky, J., Ehrlich, R., Clark, R., Le, Q.V., R\'{e}, C.,
Mirhoseini, A.: Large language monkeys: Scaling inference compute with repeated
sampling. arXiv preprint arXiv:2407.21787 (2024)

\bibitem{aggarwal2023adaptive}
Aggarwal, P., Madaan, A., Yang, Y., Mausam: Let's sample step by step:
Adaptive-consistency for efficient reasoning and coding with LLMs. In:
Proceedings of EMNLP, pp. 12375--12396 (2023)

\bibitem{wang2024tokeneconomies}
Wang, J., Jain, S., Zhang, D., Ray, B., Kumar, V., Athiwaratkun, B.:
Reasoning in token economies: Budget-aware evaluation of LLM reasoning
strategies. In: Proceedings of EMNLP, pp. 19916--19939 (2024)

\bibitem{snell2025scaling}
Snell, C., Lee, J., Xu, K., Kumar, A.: Scaling LLM test-time compute optimally
can be more effective than scaling model parameters. In: International
Conference on Learning Representations (2025)

\bibitem{damani2025hard}
Damani, M., Shenfeld, I., Peng, A., Bobu, A., Andreas, J.: Learning how hard to
think: Input-adaptive allocation of LM computation. In: International
Conference on Learning Representations (2025)

\bibitem{han2025token}
Han, T., Wang, Z., Fang, C., Zhao, S., Ma, S., Chen, Z.: Token-budget-aware
LLM reasoning. In: Findings of the Association for Computational Linguistics:
ACL, pp. 24842--24855 (2025)

\bibitem{wang2025penny}
Wang, X., Feng, S., Li, Y., Yuan, P., Zhang, Y., Tan, C., Pan, B., Hu, Y.,
Li, K.: Make every penny count: Difficulty-adaptive self-consistency for
cost-efficient reasoning. In: Findings of NAACL (2025)

\bibitem{wan2025rasc}
Wan, G., Wu, Y., Chen, J., Li, S.: Reasoning aware self-consistency:
Leveraging reasoning paths for efficient LLM sampling. In: Proceedings of
NAACL, pp. 3613--3635 (2025)

\bibitem{wu2025thought}
Wu, M., Zhou, C., Bates, S., Jaakkola, T.: Thought calibration: Efficient and
confident test-time scaling. In: Proceedings of EMNLP, pp. 14291--14305
(2025)

\bibitem{zhai2026adaptive}
Zhai, Z., Li, B., Xiao, B., Li, M., Wang, X.: Adaptive test-time compute
allocation for reasoning LLMs via constrained policy optimization. arXiv
preprint arXiv:2604.14853 (2026)

\bibitem{qwen3report}
Qwen Team: Qwen3 technical report. arXiv preprint arXiv:2505.09388 (2025)

\bibitem{cobbe2021gsm8k}
Cobbe, K., Kosaraju, V., Bavarian, M., Chen, M., Jun, H., Kaiser, L.,
Plappert, M., Tworek, J., Hilton, J., Nakano, R., Hesse, C., Schulman, J.:
Training verifiers to solve math word problems. arXiv preprint
arXiv:2110.14168 (2021)

\bibitem{hendrycks2021math}
Hendrycks, D., Burns, C., Kadavath, S., Arora, A., Basart, S., Tang, E.,
Song, D., Steinhardt, J.: Measuring mathematical problem solving with the
MATH dataset. In: NeurIPS Datasets and Benchmarks Track (2021)

\end{thebibliography}

\end{document}
